\clearpage
\SourceLeaf{495}
\pdfbookmark[1]{Nachlass V. Allgemeine Störungen der Pallas durch Jupiter. Erste Rechnung.}{pallas-V}
\fancyhead[L]{\small NACHLASS. STÖRUNGEN DER PALLAS}
\begin{center}
{\Large NACHLASS.\par}
\bigskip
{\large [V. ALLGEMEINE STÖRUNGEN DER PALLAS DURCH JUPITER. ERSTE RECHNUNG.]\par}
\medskip
[BEGONNEN 1811 AUGUST.]\par
\end{center}
\par\Needspace{9\baselineskip}\bigskip
\pdfbookmark[2]{Allgemeine Störungen, Erste Rechnung, [1]}{pallas-V-fragment1}
\begin{center}[1.]\end{center}
\begingroup
\makeatletter\tagsleft@true\let\veqno\@@leqno\makeatother
\setlength{\lineskip}{4pt}\setlength{\lineskiplimit}{4pt}
\newcommand{\VonePallasMarks}{\makebox[.16em][l]{\raisebox{1.35ex}{\rule{.07em}{.07em}}}\kern-.16em\makebox[.16em][l]{\raisebox{-.5ex}{\rule{.14em}{.035em}}}}
\noindent\hspace*{\parindent}[Zur Entwickelung der allgemeinen Theorie der Störungen der Pallas durch Jupiter wurde im Grossen und Ganzen, wie bei der Rechnung der speciellen Störungen, die in der »Exposition etc.« (Seite 449 u. 465 dieses Bandes) auseinander gesetzte Methode angewandt. Da indessen hier beabsichtigt wird, die allgemeine Theorie zu geben und nicht nur die numerischen Werthe der Störungen für einen gewissen Zeitraum zu berechnen, so werden die Gleichungen für die Störungen hier in etwas anderer Form aufgestellt werden müssen als dort. Diese sollen darum hier noch einmal zusammengestellt werden, wobei die Bezeichnungsweise dieselbe bleiben soll, wie bei den frühern Rechnungen (Seite 473 dieses Bandes).
\par
Die Coordinaten von Pallas und Jupiter \(r,v,r',w'\) berechnen sich aus den elliptischen Elementen nach bekannten Formeln; ferner zur Berechnung der gegenseitigen Bahnlage dienen die Formeln (vgl. Seite 459 dieses Bandes):
\[
\begin{aligned}
\sin\tfrac12\mathrm{J}.\sin\tfrac12(\Delta+N)&=\sin\tfrac12(i'+i).\sin\tfrac12(\AscendingNode'-\AscendingNode)\\
\sin\tfrac12\mathrm{J}.\cos\tfrac12(\Delta+N)&=\sin\tfrac12(i'-i).\cos\tfrac12(\AscendingNode'-\AscendingNode)\\
\cos\tfrac12\mathrm{J}.\sin\tfrac12(\Delta-N)&=\cos\tfrac12(i'+i).\sin\tfrac12(\AscendingNode'-\AscendingNode)\\
\cos\tfrac12\mathrm{J}.\cos\tfrac12(\Delta-N)&=\cos\tfrac12(i'-i).\cos\tfrac12(\AscendingNode'-\AscendingNode),
\end{aligned}
\tag*{1)}
\]
aus denen sich
\begin{center}
\begin{tabular}{@{}r@{\,}p{.82\linewidth}@{}}
\(\mathrm{J}={}\)&Neigung beider Bahnebenen,\\
\(\AscendingNode+\Delta={}\)&Länge des Knotens der Jupitersbahn auf der \VonePallasMarks Pallasbahn, gezählt auf der letztern,\\
\(\AscendingNode'+N={}\)&»\quad »\quad »\quad »\quad »\quad »\quad »\quad »\quad »\quad »\quad »\quad erstern
\end{tabular}
\end{center}
ergeben. Weiter hat man zur Berechnung von \(v'\) und \(\beta'\) die Gleichungen
\[
\begin{aligned}
\operatorname{tang}(v'-\AscendingNode-\Delta)&=\cos\mathrm{J}.\operatorname{tang}(w'-\AscendingNode'-N)\\
\operatorname{tang}\beta'&=\operatorname{tang}\mathrm{J}.\sin(v'-\AscendingNode-\Delta)\\
\sin\beta'&=\sin\mathrm{J}.\sin(w'-\AscendingNode'-N),
\end{aligned}
\tag*{2)}
\]
\SourceLeaf{496}von denen eine der beiden letztern zur Controlle dienen kann. Ferner findet sich
\[
\begin{aligned}
\xi&=r'\cos\beta'\cos(v-v')&\quad\rho\rho&=r'r'+rr-2r\xi\\[.4em]
\eta&=r'\cos\beta'\sin(v-v')&\Delta&=\frac{1}{\rho^3}-\frac{1}{r'^3}\\[.4em]
\zeta&=r'\sin\beta'&&\\[.4em]
T&=\frac{\alpha r}{\rho^3}-\alpha\Delta\xi&&\\[.4em]
V&=\alpha r\Delta\eta&\alpha&=\frac{\mu a}{\cos\varphi}\mathpunct{\raisebox{.45ex}{.}}\\[.4em]
W&=\alpha r\Delta\zeta&&
\end{aligned}
\tag*{3)}
\]
Der Werth von \(\alpha\) ist noch mit der Zahl \(206265''\) zu multipliciren, womit man alle Grössen ausgedrückt in Bogensecunden erhält.
\par
Nach den Gleichungen 3, Seite 474, gelten dann für die Störungen der Neigung und der Knotenlänge die Relationen:
\[
\begin{aligned}
\frac{\mathrm{d}i}{n\mathrm{d}t}&=\cos(v-\AscendingNode).W\\[.5em]
\frac{\mathrm{d}\AscendingNode}{n\mathrm{d}t}&=\frac{1}{\sin i}\sin(v-\AscendingNode).W.
\end{aligned}
\tag*{4)}
\]
\par
Für die Störungen des Perihels ergibt eine leichte Transformation der genannten Gleichungen:
\[
\begin{aligned}
\frac{\mathrm{d}\tilde\omega}{n\mathrm{d}t}
&=(1-\cos i)\frac{\mathrm{d}\AscendingNode}{n\mathrm{d}t}
+\frac{p}{e}\cos(v-\tilde\omega).T\\
&\qquad-\frac{1}{e}\{2+e\cos(v-\tilde\omega)\}\sin(v-\tilde\omega).V.
\end{aligned}
\tag*{5)}
\]
\par
Für den halben Parameter \(p\), der an Stelle der Excentricität eingeführt werden soll, ergibt Gleichung 63 der Exposition mit Berücksichtigung der veränderten Bedeutung von \(V\):
\[
\frac{\mathrm{d}p}{n\mathrm{d}t}=-2a\cos\varphi^2.V,
\]
also
\[
\frac{\mathrm{d}\log p}{n\mathrm{d}t}=-2V.
\tag*{6)}
\]
\par
Des Weitern ergeben die Gleichungen 3, Seite 474,
\[
\begin{aligned}
\frac{\mathrm{d}\log\operatorname{hyp} n}{n\mathrm{d}t}
&=3ae\sin(v-\tilde\omega).T+\frac{3aa\cos\varphi^2}{rr}V\\[.7em]
\frac{\mathrm{d}\varepsilon}{n\mathrm{d}t}
&=\operatorname{tang}\tfrac12 i\sin(v-\AscendingNode).W\\
&\quad+\{2r\cos\varphi+p\operatorname{tang}\tfrac12\varphi\cos(v-\tilde\omega)\}T\\
&\quad-\operatorname{tang}\tfrac12\varphi(2+e\cos(v-\tilde\omega))\sin(v-\tilde\omega).V\\[.5em]
&=(1-\cos\varphi)\left(\frac{\mathrm{d}\tilde\omega}{n\mathrm{d}t}
 -(1-\cos i)\frac{\mathrm{d}\AscendingNode}{n\mathrm{d}t}\right)\\
&\quad+(1-\cos i)\frac{\mathrm{d}\AscendingNode}{n\mathrm{d}t}
 +2r\cos\varphi.T.
\end{aligned}
\tag*{7)}
\]
\par
Legt man nun der ganzen Rechnung sowohl für Pallas wie für Jupiter ein elliptisches Elementensystem zum Grunde, so sind offenbar die Grössen \(T,V,W\) ebenso wie die Differentialquotienten der Elemente periodische Functionen der mittlern Anomalien beider Planeten, also der Grössen \(M\) und \(M'\), und lassen sich durch trigonometrische Reihen darstellen nach Vielfachen dieser beiden Argumente. Man kann darum hier die in der Theoria interpolationis methodo nova tractata (Band III, Seite 265) gemachten Untersuchungen benutzen und die Coefficienten dieser Reihen bestimmen, indem man die numerischen Werthe von \(T,V,W\) für eine Reihe von Werthen von \(M\) und \(M'\) berechnet, die je eine ganze Periode umfassen.
\par
\SourceLeaf{497}Zu diesem Zwecke sollen die in der genannten Abhandlung abgeleiteten Formeln in den folgenden Artikeln noch einer Umformung unterworfen werden, durch welche sie für die Lösung der vorliegenden Aufgabe möglichst einfach gestaltet werden.]
\par
\endgroup
